\scanpage{6}

\deleted{$F'(T)$ l'ensemble analogue, mais où dans la structure de $C$ on se donne de plus une trivialisation \deleted{\unsure{par}} de \deleted{$\unsure{\Omega^1|\underline{\Omega}}$} l'image directe de $(\Omega^1_{C/T})^{\otimes\unsure{h|k}}$ ($h\geq 3$ étant choisi). Soit
\[
\unsure{\nu|d|v}=2h(g-1)-(g-1)=(2h-1)(g-1)
\]}

\deleted{Proposition} \inserted{\uline{Théorème}} \unsure{Les|Le} foncteurs $F(T)$ \deleted{et $\unsure{F'(T)|F(T)}$ sont} \inserted{est} représentable, \deleted{et} en un $M$ quasi projectif sur $S$ ; le morphisme \deleted{c'est un fibré principal homogène en $F(T)$ \unsure{est} \inserted{représenté \unsure{par}} \uline{un sous-schéma} (\unsure{au} \unsure{quasi} \unsure{non fermé}) de $S$, et $F'(T)$ est un fibré principal homogène sur $F(T)$, de groupe $\underline{\underline{\mathrm{GL}}}(\unsure{\nu})_S$} \inserted{structural $M\to S$ est radiciel et \unsure{a une|à une} \unsure{image} \deleted{\unsure{relativement|rel.}} fermée.}
% slanted
\margin{\unsure{rerédiger|rédiger} proprement}

Plus précisément : le morphisme structural \deleted{\unsure{est}} \inserted{\unsure{induit}} \deleted{\unsure{un}} \inserted{un} morphisme \deleted{\unsure{propre et quasi-fini radiciel}} \inserted{fini} sur un \unsure{sous-préschéma|sous-schéma} \unsure{fermé} \unsure{(ou ouvert)|(ouvert)} de $S$ ; \unsure{ce|le} morphisme est \deleted{un isomorphisme} \inserted{non ramifié} en \unsure{tt pt géom.}
% slanted
\margin{\uline{\unsure{ses pts|ex :} correspondant à une courbe non hyperelliptique}, \deleted{\illegible{3}} \uline{sa fibre} \deleted{\illegible{1}} au dessus d'un pt hyperelliptique \unsure{a un seul point} \deleted{\illegible{2}} \uline{géom.}, \unsure{isom. à} $\unsure{\mathrm{Spec}\,K[t]/t^2}$.}

\uline{Démonstration} \deleted{\unsure{1)} Commençons par représenter $F'(T)$, en commençant par remplacer le schéma $M$ sur $S$ Soit $T$ \unsure{faisons} le changement de base $T$} Considérons le foncteur $F'(T)$ analogue à $F(T)$, où \unsure{l'on} \deleted{nouveau foncteur représenté} met sur $C$ une structure supplémentaire : une section de $J^1_{C/S}$ sur $T$. \deleted{\unsure{D'après le numéro précédent|D'après les remarques générales} Il y a donc une \unsure{projection|morphisme} naturelle $F'(\deleted{\unsure{T'|T}}T)\to F(T)$, le produit fibré $\unsure{F'(T)\times_F|F(T)\times_F}$} $F'\times_F F'\mathrel{\unsure{\rightrightarrows|=}}F'$ étant \unsure{un} d'ailleurs, \unsure{si|d$^\circ$} $F'$ est représentable \inserted{par $M'$}, identifiable à \deleted{\unsure{un schéma}} $A\times M'$, \unsure{opérant sur|opérant en} $M'$ par des opérations de groupe.
% slanted
\margin{\deleted{\uline{\unsure{développer}} \uline{\unsure{en détail}}}}

% archive sheet number: 5
